\section{Posterior pricing and the integration correction}\label{sec:short_method}

\subsection{Historical inference and conditional valuation}
Historical log returns follow the physical-measure benchmark
\begin{equation}
 R_i\sim\mathcal N\!\left[
 \left(\mu-\frac12\sigma^2\right)\Delta t,
 \sigma^2\Delta t\right],
 \label{eq:short_returns}
\end{equation}
where $\mu$ is annualized drift, $\sigma>0$ is annualized volatility, and $\Delta t$ is the observation interval in years. The baseline priors are
\begin{equation}
 \mu\sim\mathcal N(0,1),\qquad
 \sigma\sim\operatorname{InvGamma}(2,0.1).
 \label{eq:short_priors}
\end{equation}
The inverse-gamma prior is placed directly on volatility, with shape $2$ and scale $0.1$. Empirical sampling uses random-walk Metropolis in $(\mu,\log\sigma)$, including the transformation Jacobian. For the synthetic experiment, $\mu$ is marginalized analytically conditional on $\sigma$ and the remaining posterior is evaluated by deterministic quadrature.

Under the risk-neutral measure, futures maturity $u$ follows
\begin{equation}
 dF_t^u=\sigma F_t^u\,dW_t^Q.
 \label{eq:short_q}
\end{equation}
The same Brownian motion drives all maturities. This common-factor assumption fixes the cross-maturity dependence used to value the average. Historical drift does not enter the valuation. Using the historically inferred $\sigma$ in Equation~\eqref{eq:short_q} defines the historical-volatility benchmark; its empirical adequacy is evaluated separately from the effect of integrating its posterior.

For fixing dates $t_1,\ldots,t_M$, the settlement average is
\begin{equation}
 A_T=\frac1M\sum_{j=1}^M F^{(1)}_{t_j},
 \label{eq:short_average}
\end{equation}
where $F^{(1)}_{t_j}$ is the first-nearby CL futures settlement applicable on date $t_j$. At valuation date $t$, realized fixings enter deterministically and unresolved fixings use their applicable futures maturities initialized from the current CL curve. The conditional call value is
\begin{equation}
 C_t^Q(\sigma)=B(t,T)\,
 \mathbb E^Q[(A_T-K)^+\mid\mathcal F_t;\sigma],
 \label{eq:short_call_value}
\end{equation}
with strike $K$, discount factor $B(t,T)$, and valuation information $\mathcal F_t$. Put values replace $(A_T-K)^+$ with $(K-A_T)^+$. Below, $C^Q(\sigma)$ suppresses the fixed contract state and applies to either option type.

Historical PI and PM are evaluated by Monte Carlo. Curran's conditioning approximation \citep{curran1994} supplies deterministic prices for scalar historical and option-implied inputs. Numerical checks compare Monte Carlo, independently scrambled Sobol pricing, and Curran calculations. Pricing approximation and posterior integration are distinct effects; the reported pricing method is identified in each empirical table.

\subsection{Point-price correction versus price dispersion}
Writing $\delta=\sigma-\bar\sigma$, a Taylor expansion gives
\begin{align}
 C^Q(\sigma)={}&C^Q(\bar\sigma)+C^{Q\prime}(\bar\sigma)\delta\notag\\
 &+\tfrac12 C^{Q\prime\prime}(\bar\sigma)\delta^2+R_3(\sigma).
 \label{eq:short_expansion}
\end{align}
Because $\mathbb E[\delta\mid\mathcal D]=0$, averaging this expression yields Equation~\eqref{eq:short_taylor} plus $\mathbb E[R_3\mid\mathcal D]$. When the third derivative is bounded in absolute value by $M_3$ over the relevant range,
\begin{equation}
 |\mathbb E[R_3\mid\mathcal D]|
 \leq\frac{M_3}{6}\mathbb E[|\delta|^3\mid\mathcal D].
 \label{eq:short_remainder}
\end{equation}
The approximation is useful when the posterior is sufficiently concentrated and higher-order curvature is controlled. The integration correction can have either sign; positive volatility curvature is not assumed.

The same posterior induces conditional values $C^Q(\sigma)$ with dispersion
\begin{equation}
 \operatorname{SD}(C^Q(\sigma)\mid\mathcal D)
 \approx |C^{Q\prime}(\bar\sigma)|
 \sqrt{\operatorname{Var}(\sigma\mid\mathcal D)}.
 \label{eq:short_delta}
\end{equation}
With bounded derivatives, the integration correction is locally second order in posterior standard deviation, whereas price dispersion has a first-order term when volatility sensitivity is nonzero. Thus PI--PM measures a shift in the point value, not the full uncertainty around it. Posterior model-price intervals quantify uncertainty in this conditional valuation; they are not predictive intervals for future market settlements.

\subsection{Mapping information and contract state}
The synthetic experiment contains 175,000 independent histories: 5,000 for each combination of $n\in\{21,42,63,126,252,504,1260\}$ and $\sigma_0\in\{0.10,0.20,0.35,0.50,0.80\}$. The contract layer crosses four horizons from 21 to 252 days, seven moneyness values from 0.60 to 1.50, and six partial-fixing states. The forward level is normalized to \$100 per barrel. For each history and state, the experiment compares the PI--PM correction with posterior variance times local volatility curvature. The experiment measures the approximation's performance and the magnitude of the correction across these states; the expansion itself is a standard analytical device.
